\scanpage{121}\mspage{61}

Cela se voit facilement si l'ordre des points de $P_0$ est impair, car alors $G_0^+$ se décrit intrinsèquement \deleted{comme} en termes de $G_{\unsure{0}}$, par le fait que c'est le \uline{seul} sous-groupe d'indice 2. Si par contre $n$ est pair, ce n'est pas chose évidente -- il y a \uline{trois} sous-groupes d'indice 2 possibles \inserted{a priori}, car $(\unsure{\mathbb G_{\infty\,3}|G_{\infty\,3}})_{\mathrm{ab}}\simeq\mathbb F_2\times\mathbb F_2$. Mais il n'y en a qu'un seul qui contienne $\rho_1^r$, donc la question est si $\rho_1^r$ (relatif à un repère $r\deleted{\unsure{\}}}=\{s_0,t_0,u_0\}$ de $P_0\}$ conserve l'orientation \deleted{dans} $P_\lambda$\ldots

% slanted
\margin{il faudra revenir là-dessus\ldots}

Je m'intéresse maintenant à l'ensemble des structures $P_\lambda$, \inserted{$\lambda\in Z'$.} \deleted{Notons d'abord} Les structures polyédrales triangulées qui satisfont aux axiomes $(T_0\,1)$ à $(T_0\,6)$ seront appelées (à défaut de mieux, et provisoirement) ``spéciales''. Je dirais que deux telles structures, \inserted{définies par des ens. d'arêtes $\subset\unsure{\mathbb P_2|\mathcal P_2|\mathfrak P_2}(S)$} $\deleted{\unsure{F_0}}A_0,\deleted{\unsure{F_1}}A_1$ sur un même ensemble de sommets \inserted{$S$} sont \uline{équivalentes}, il existe un $\lambda\in\unsure{\mathfrak S_S|G_S}$, commutant à $\operatorname{Aut}(S,F_0)$, de telle façon que l'on ait
% archive sheet number: 120
